% Luc Destombe --- licence CC BY-NC-SA 4.0
% https://creativecommons.org/licenses/by-nc-sa/4.0/deed.fr
% Fichier autonome : compiler avec pdflatex, deux fois (signets, nombre de pages).
\documentclass[11pt,a4paper]{article}
\makeatletter
\newif\ifeleve
\elevefalse

\newif\iflycee@palatino
\lycee@palatinotrue

\RequirePackage[utf8]{inputenc}
\RequirePackage[T1]{fontenc}
\RequirePackage[french]{babel}
\RequirePackage{etoolbox}

\newcommand{\ShorthandsOff}{\@ifundefined{shorthandoff}{}{\shorthandoff{;:!?}}}
\newcommand{\ShorthandsOn}{\@ifundefined{shorthandon}{}{\shorthandon{;:!?}}}
\ShorthandsOff
\AtBeginDocument{\ShorthandsOn}
\AtBeginEnvironment{tikzpicture}{\ShorthandsOff}

\RequirePackage{geometry}
\RequirePackage{fancyhdr}
\RequirePackage{lastpage}
\RequirePackage{multicol}
\RequirePackage{array}
\RequirePackage{enumitem}
\RequirePackage{xcolor}
\RequirePackage{needspace}

\RequirePackage{amsmath,amssymb}
\RequirePackage{mathpazo}
\DeclareMathSizes{10.95}{10.95}{8}{5.5}
\begingroup\catcode`\_=\active \catcode`\^=\active
\gdef_#1{\sb{\scriptscriptstyle #1}}
\gdef^#1{\sp{\scriptscriptstyle\lycee@moinsepais #1}}
\endgroup
\newcommand\lycee@moins{\mathbin{%
  \setbox\z@\hbox{$\m@th\scriptscriptstyle\lycee@moinsorig$}%
  \dimen@=.55\wd\z@ \advance\dimen@-.5pt
  \hbox to.75\wd\z@{\hss$\m@th\scriptscriptstyle\vcenter{\hbox{%
    \tikz\draw[line width=.5pt,line cap=round](0,0)--(\the\dimen@,0);}}$\hss}}}
\begingroup\lccode`\~=`\- \lowercase{\endgroup\let~}\lycee@moins
\newcommand\lycee@moinsepais{\mathcode`\-="8000 }
\AtBeginDocument{\mathchardef\lycee@moinsorig=\mathcode`\-\relax}
\AtBeginDocument{\catcode`\_=12 \mathcode`\_="8000
  \catcode`\^=12 \mathcode`\^="8000 }
\newcommand\lycee@exposants{%
  \fontdimen13\textfont2=.48\fontdimen6\textfont2
  \fontdimen14\textfont2=.48\fontdimen6\textfont2
  \fontdimen15\textfont2=.48\fontdimen6\textfont2
  \fontdimen13\scriptfont2=.48\fontdimen6\scriptfont2
  \fontdimen14\scriptfont2=.48\fontdimen6\scriptfont2
  \fontdimen15\scriptfont2=.48\fontdimen6\scriptfont2}
\every@math@size\expandafter{\the\every@math@size\lycee@exposants}
\newlength{\retraitcalculs}
\setlength{\retraitcalculs}{2em}

\RequirePackage{tikz}
\usetikzlibrary{arrows.meta,calc,babel,shapes.misc}
\RequirePackage[most]{tcolorbox}
\tcbuselibrary{skins,breakable}

\definecolor{coulDef}{HTML}{1F4E79}
\definecolor{coulProp}{HTML}{7030A0}
\definecolor{coulRem}{HTML}{C55A11}
\definecolor{coulEx}{HTML}{2E7D32}
\definecolor{coulExo}{HTML}{B03A2E}
\definecolor{coulPart}{HTML}{1F4E79}
\definecolor{coulMeth}{HTML}{1B6E4F}
\definecolor{coulCourbe}{HTML}{0B5ED7}
\definecolor{coulCourbeB}{HTML}{C00000}
\definecolor{coulCourbeC}{HTML}{2E7D32}
\definecolor{coulGrille}{HTML}{8C8C8C}
\definecolor{coulRep}{HTML}{1B6E4F}
\definecolor{coulBar}{HTML}{2E7D32}

\newcommand{\R}{\mathbb{R}}
\newcommand{\Rs}{\mathbb{R}^{*}}
\renewcommand{\le}{\leqslant}
\renewcommand{\ge}{\geqslant}
\renewcommand{\approx}{\simeq}
\newcommand{\vect}[1]{\overrightarrow{#1}}
\newcommand{\norme}[1]{\left\lVert #1 \right\rVert}
\newcommand{\intervalle}[2]{\left[#1\,;#2\right]}
\RequirePackage{cancel}
\renewcommand{\CancelColor}{\color{coulExo}}
\newcommand{\fois}[1]{\times{\color{coulExo}#1}}
\newcommand{\barre}[1]{\cancel{#1}}

\tikzset{grille/.style={coulGrille,line width=0.4pt}}
\newcommand{\quadrillage}[4]{\draw[grille,step=1] (#1,#2) grid (#3,#4);}
\tikzset{
  croix/.style={cross out,draw,line width=0.6pt,inner sep=0pt,minimum size=4pt},
  croixdroite/.style={croix,rotate=45,line width=0.8pt},
}
\newcommand{\pt}[3]{\path (#1) node[croix]{} node[#2,font=\small,inner sep=2.5pt]{$#3$};}

\newcommand{\lycee@repere}[4]{%
  \draw[grille,step=1] (#1,#2) grid (#3,#4);
  \draw[-{Stealth[length=2mm]},thick] (#1,0)--({#3+0.4},0) node[below left=-1pt]{$x$};
  \draw[-{Stealth[length=2mm]},thick] (0,#2)--(0,{#4+0.4}) node[below left=-1pt]{$y$};
  \node[below left,font=\scriptsize,inner sep=1.5pt] at (0,0) {$0$};}
\newcommand{\lycee@gradx}[1]{\foreach \k/\l in {#1}{%
  \draw (\k,0.07)--(\k,-0.07) node[below,font=\scriptsize,inner sep=1.5pt]{$\l$};}}
\newcommand{\lycee@grady}[1]{\foreach \k/\l in {#1}{%
  \draw (0.07,\k)--(-0.07,\k) node[left,font=\scriptsize,inner sep=1.5pt]{$\l$};}}
\newcommand{\lycee@intff}[2]{\left[#1\,;#2\right]}
\newcommand{\lycee@intfo}[2]{\left[#1\,;#2\right[}
\newcommand{\lycee@intof}[2]{\left]#1\,;#2\right]}
\newcommand{\lycee@intoo}[2]{\left]#1\,;#2\right[}
\newcommand{\lycee@ens}[1]{\left\{#1\right\}}
\AtBeginDocument{%
  \providecommand{\repere}{\lycee@repere}%
  \providecommand{\gradx}{\lycee@gradx}%
  \providecommand{\grady}{\lycee@grady}%
  \providecommand{\Cf}{\mathcal{C}\sb{\scriptscriptstyle f}}%
  \providecommand{\Cg}{\mathcal{C}\sb{\scriptscriptstyle g}}%
  \providecommand{\intff}{\lycee@intff}%
  \providecommand{\intfo}{\lycee@intfo}%
  \providecommand{\intof}{\lycee@intof}%
  \providecommand{\intoo}{\lycee@intoo}%
  \providecommand{\ens}{\lycee@ens}}

\newlist{questions}{enumerate}{3}
\setlist[questions,1]{label=\arabic*\textdegree),leftmargin=*,itemsep=2pt,topsep=3pt}
\setlist[questions,2]{label=\alph*),leftmargin=*,itemsep=1pt,topsep=2pt}
\setlist[questions,3]{label=\roman*),leftmargin=*,itemsep=1pt,topsep=1pt}
\setlist[itemize]{label=\textbullet,leftmargin=*,itemsep=2pt,topsep=3pt}

\newcommand{\besoinplace}[1]{\par\penalty\@M\begingroup
  \dimen@=#1\relax\dimen@ii=\pagegoal\advance\dimen@ii-\pagetotal
  \ifdim\dimen@>\dimen@ii\ifdim\dimen@ii>\z@\vfil\fi\break\fi\endgroup}

\newcommand{\lycee@pied}{\thepage/\pageref{LastPage}}
\newcommand{\entete}[2]{%
  \pagestyle{fancy}\fancyhf{}%
  \fancyhead[L]{\footnotesize\color{coulPart}\textbf{#1}}%
  \fancyhead[R]{\footnotesize\color{coulPart}\textbf{#2}}%
  \fancyfoot[C]{\footnotesize\lycee@pied}%
  \renewcommand{\headrulewidth}{0.4pt}%
  \renewcommand{\headrule}{\hbox to\headwidth{%
    \color{coulPart!40}\leaders\hrule height \headrulewidth\hfill}}}

\RequirePackage{ccicons}
\newcommand{\lycee@licence}{{\scriptsize\color{gray}%
  \copyright\ Luc Destombe --- \ccbyncsa\ CC BY-NC-SA 4.0}}
\newif\iflycee@licence \lycee@licencetrue
\newcommand{\sanslicence}{\lycee@licencefalse}
\AtBeginDocument{\iflycee@licence\AddToHookNext{shipout/foreground}{%
  \put(0,0){\rlap{\kern\dimexpr1in+\hoffset+\oddsidemargin+\textwidth\relax
    \llap{\raisebox{-\dimexpr1in+\voffset+\topmargin+\headheight+\headsep
      +\textheight+\footskip\relax}{\lycee@licence}}}}}\fi}

\newcommand{\formulesserrees}{%
  \setlength{\abovedisplayskip}{3pt}\setlength{\belowdisplayskip}{3pt}%
  \setlength{\abovedisplayshortskip}{2pt}\setlength{\belowdisplayshortskip}{3pt}}

\setlength{\parindent}{0pt}

\RequirePackage[implicit=false,hidelinks,pdfencoding=auto,psdextra,bookmarksopen,
  bookmarksopenlevel=1,pdfcreator={},pdfproducer={}]{hyperref}
\RequirePackage{bookmark}
\hypersetup{pdfauthor={Luc Destombe},pdfkeywords={CC BY-NC-SA 4.0}}
\newcounter{lycee@signet}
\newcommand{\signet}[3][]{\stepcounter{lycee@signet}%
  \edef\lycee@dest{\if\relax\detokenize{#1}\relax
    signet.\arabic{lycee@signet}\else#1\fi}%
  \hypertarget{\lycee@dest}{}\bookmark[level=#2,dest=\lycee@dest]{#3}}
\pdfstringdefDisableCommands{%
  \def\R{R}\def\vect#1{#1}\def\Cf{Cf}\def\Cg{Cg}\def\fois#1{×#1}%
  \def\barre#1{#1}\def\intervalle#1#2{[#1;#2]}\def\intff#1#2{[#1;#2]}%
  \def\textsuperscript#1{#1}\def\dfrac#1#2{#1/#2}\def\frac#1#2{#1/#2}%
  \def\vec#1{#1⃗}}%
\geometry{top=2cm,bottom=1cm,left=1cm,right=1cm,headheight=14pt,footskip=0.6cm}

\RequirePackage{pgfplots}
\pgfplotsset{compat=1.18}
\RequirePackage{tkz-tab}

\newcounter{partiecours}
\renewcommand{\thepartiecours}{\Roman{partiecours}}
\newcommand{\partiecours}[1]{%
  \refstepcounter{partiecours}%
  \Needspace*{8\baselineskip}%
  \bigskip
  \noindent\begin{tcolorbox}[enhanced,colback=coulPart,colframe=coulPart,
      boxrule=0pt,arc=2pt,left=8pt,right=8pt,top=4pt,bottom=4pt,
      before skip=6pt,after skip=10pt]
    \color{white}\bfseries\large \thepartiecours{} --- #1%
    \signet[partie-\arabic{partiecours}]{1}{\thepartiecours{} --- #1}
  \end{tcolorbox}%
  \addcontentsline{toc}{section}{\thepartiecours{} --- #1}%
}

\newtcolorbox{boitecours}[3][]{%
  enhanced,breakable,
  colback=white,colframe=#2,coltitle=white,
  fonttitle=\bfseries,
  attach boxed title to top left={xshift=8pt,yshift=-\tcboxedtitleheight/2},
  boxed title style={colback=#2,arc=2pt,boxrule=0pt},
  boxrule=0.9pt,arc=2pt,
  left=8pt,right=8pt,top=8pt,bottom=6pt,
  title={#3},#1}

\newcounter{methode}
\newenvironment{definitionbox}[1][Définition]%
  {\begin{boitecours}[phantom={\signet{2}{#1}}]{coulDef}{#1}}{\end{boitecours}}
\newenvironment{proprietebox}[1][Propriété]%
  {\begin{boitecours}[phantom={\signet{2}{#1}}]{coulProp}{#1}}{\end{boitecours}}
\newenvironment{methodebox}[1][Méthode]{\stepcounter{methode}%
  \begin{boitecours}[phantom={\signet[methode-\arabic{methode}]{2}{#1}}]{coulMeth}{#1}}%
  {\end{boitecours}}

\newcommand{\etape}[1]{\par\smallskip\textbf{\color{coulMeth}#1}\par\nobreak\smallskip}
\newcommand{\enonce}[1]{\emph{#1}\par\smallskip}
\newcommand{\reponse}[1]{\par\smallskip
  {\footnotesize\itshape\color{black!55}Réponses : #1}}
\newcommand{\demo}[1]{\textbf{\color{coulMeth}Démonstration.} #1}
\newcommand{\afaire}[1]{%
  \par\smallskip
  \noindent{\small\color{coulExo}$\blacktriangleright$\ \textbf{À faire :}
    fiche d'exercices du chapitre, #1.}\par\smallskip}

\newenvironment{calculs}{\setlength{\jot}{5pt}\@fleqntrue\@mathmargin\retraitcalculs
  \setlength{\abovedisplayskip}{4pt}\setlength{\belowdisplayskip}{4pt}%
  \setlength{\abovedisplayshortskip}{4pt}\setlength{\belowdisplayshortskip}{4pt}%
  \csname alignat*\endcsname{2}}%
  {\csname endalignat*\endcsname}
\newcommand{\rem}[1]{\text{\small\itshape\color{black!60}#1}}
\newcommand{\explique}[2]{\par\smallskip\noindent
  \begin{minipage}[t]{0.57\linewidth}\raggedright #1\end{minipage}\hfill
  \begin{minipage}[t]{0.40\linewidth}\raggedright\leavevmode
    {\small\itshape\color{black!60}#2}\end{minipage}\par}
\newcommand{\cas}[1]{\par\smallskip\noindent\textbf{\color{coulExo}#1}\nobreak}
\newcommand{\calc}[1]{\par\smallskip\textbf{\color{coulExo}#1}\par\nobreak}

\newtcolorbox{remarquebox}[1][Remarque]{%
  enhanced,breakable,colback=white,colframe=coulRem,
  boxrule=0pt,leftrule=2.5pt,arc=0pt,outer arc=0pt,
  left=8pt,right=6pt,top=5pt,bottom=5pt,
  before upper={\textbf{\color{coulRem}#1}\par\smallskip}}

\newtcolorbox{exemplebox}[1][Exemple]{%
  enhanced,breakable,colback=white,colframe=coulEx,
  boxrule=0pt,leftrule=2.5pt,arc=0pt,outer arc=0pt,
  left=8pt,right=6pt,top=5pt,bottom=5pt,
  before upper={\textbf{\color{coulEx}#1}\par\smallskip}}

\pgfplotsset{
  repere/.style={
    axis lines=middle,
    axis line style={-{Stealth[length=2.4mm]},thick},
    every axis x label/.style={at={(ticklabel* cs:1.0)},anchor=west},
    every axis y label/.style={at={(ticklabel* cs:1.0)},anchor=south},
    label style={font=\footnotesize},
    tick label style={font=\scriptsize},
    grid=both,
    major grid style={grille},
    minor tick num=0,
    tick align=inside,
    clip mode=individual,
  },
}
\tikzset{
  courbe/.style={coulCourbe,line width=1.1pt,smooth},
  courbeB/.style={coulCourbeB,line width=1.1pt,smooth},
  courbeC/.style={coulCourbeC,line width=1.1pt,smooth},
}

\newlist{etapes}{enumerate}{1}
\setlist[etapes]{label=\textbf{\arabic*.},leftmargin=*,itemsep=1pt,topsep=2pt}

\newlist{expressions}{enumerate}{1}
\setlist[expressions]{label={\Alph*\ $=$},leftmargin=*,itemsep=3pt,topsep=3pt}

\newcounter{exo}
\renewcommand{\theexo}{\ifnum\value{exo}<10 0\fi\arabic{exo}}
\newtcolorbox{exercice}[1][]{%
  enhanced,breakable,
  colback=white,colframe=coulExo,
  boxrule=0.9pt,arc=2pt,
  left=8pt,right=8pt,top=8pt,bottom=6pt,
  coltitle=white,fonttitle=\bfseries,
  attach boxed title to top left={xshift=8pt,yshift=-\tcboxedtitleheight/2},
  boxed title style={colback=coulExo,arc=2pt,boxrule=0pt},
  phantom={\signet{2}{Exercice \arabic{exo}}},
  title={Exercice \theexo\ifx\relax#1\relax\else{} --- #1\fi}}
\newenvironment{exo}[1][\relax]{\refstepcounter{exo}\begin{exercice}[#1]}{\end{exercice}}

  \newenvironment{correction}{\par}{\par}
  \newcommand{\autableau}{}
  \newcommand{\etapeexemples}{\etape{Exemples corrigés}}
  \newcommand{\etapetableau}[1]{\etape{#1}}
  \newenvironment{exempletableau}[1][Exemple]%
    {\begin{exemplebox}[#1]}{\end{exemplebox}}

\RequirePackage{comment}
\excludecomment{correctionavj}

\newcommand{\coursEntete}{}
\newcommand{\coursNiveau}{}
\newcommand{\coursTitre}{}
\newcommand{\coursSurTitre}{}
\newcommand{\coursSousTitre}{}

\newcommand{\enteteCours}[1]{\renewcommand{\coursEntete}{#1}}
\newcommand{\chapitrecours}[2]{\enteteCours{Chapitre #1 --- #2}}
\newcommand{\niveaucours}[1]{\renewcommand{\coursNiveau}{#1}}
\newcommand{\titrecours}[3][]{%
  \renewcommand{\coursTitre}{#2}%
  \renewcommand{\coursSurTitre}{#1}%
  \renewcommand{\coursSousTitre}{#3}}

\pagestyle{fancy}
\fancyhf{}
\fancyhead[L]{\footnotesize\color{coulPart}\textbf{\coursEntete}}
\fancyhead[R]{\footnotesize\color{coulPart}\textbf{\coursNiveau}}
\fancyfoot[C]{\footnotesize\thepage}
\renewcommand{\headrulewidth}{0.4pt}
\renewcommand{\headrule}{\hbox to\headwidth{\color{coulPart!40}\leaders\hrule height \headrulewidth\hfill}}

\setlength{\parskip}{2pt}

\AtBeginDocument{%
  \begin{center}
    \begin{tcolorbox}[enhanced,width=0.72\linewidth,colback=white,colframe=coulPart,
        boxrule=1.6pt,arc=3pt,halign=center,top=10pt,bottom=10pt]
      \ifx\coursSurTitre\empty
        {\Huge\bfseries\color{coulPart} \coursTitre}\\[4pt]
      \else
        {\Huge\bfseries\color{coulPart} \coursTitre}\\[3pt]
        {\Large\bfseries\color{coulPart} \coursSurTitre}\\[5pt]
      \fi
      {\normalsize \coursSousTitre}%
      
    \end{tcolorbox}
  \end{center}%
}
\makeatother
\usepackage{colortbl}

\chapitrecours{2}{Intervalles et fonctions}
\niveaucours{Seconde}
\titrecours{INTERVALLES ET FONCTIONS}{Cours --- Seconde}

\begin{document}

\ShorthandsOff
\definecolor{coulInt}{HTML}{D62828}
\definecolor{coulIntB}{HTML}{1565C0}
\providecommand{\axegrad}[2]{%
  \draw[-{Stealth[length=2mm]}] ({#1-0.5},0)--({#2+0.5},0);
  \foreach \k in {#1,...,#2} \draw (\k,0.07)--(\k,-0.07);}
\ProvideDocumentCommand{\morceau}{O{coulInt} O{0} m m}{%
  \draw[#1,line width=1.8pt] (#3,#2)--(#4,#2);}
\ProvideDocumentCommand{\crochet}{O{coulInt} O{0} m m}{%
  \ifx[#4%
    \draw[#1,line width=1.2pt] ($(#3,#2)+(2.6pt,5pt)$) -- ++(-2.6pt,0) -- ++(0,-10pt) -- ++(2.6pt,0);
  \else
    \draw[#1,line width=1.2pt] ($(#3,#2)+(-2.6pt,5pt)$) -- ++(2.6pt,0) -- ++(0,-10pt) -- ++(-2.6pt,0);
  \fi}
\ProvideDocumentCommand{\bornelab}{O{0} m m}{%
  \node[below=6.5pt,font=\scriptsize,inner sep=1pt] at (#2,#1) {$#3$};}
\providecommand{\surligne}[2]{\fill[yellow!45] (#1,-0.55) rectangle (#2,0.55);}
\providecommand{\mini}[1]{\begin{tikzpicture}[x=0.42cm,baseline=-3pt]
  \draw[-{Stealth[length=1.6mm]}] (-2.6,0)--(5.6,0);
  \foreach \k in {-2,...,5} \draw (\k,0.06)--(\k,-0.06);
  #1\end{tikzpicture}}
\providecommand{\resultat}[1]{\colorbox{coulMeth!12}{$#1$}}
\ShorthandsOn

\begin{remarquebox}[Comment utiliser ce cours]
Chaque partie commence par ce qu'il faut \textbf{savoir} (définitions, propriétés), puis
vient ce qu'il faut \textbf{savoir faire} : une méthode, avec des
\textbf{exemples corrigés} faits en classe, puis des exercices
\og\textbf{À vous de jouer}\fg{} à chercher seul.

\end{remarquebox}

\partiecours{Droite numérique, intervalles fermés}

\begin{definitionbox}[Définition 1 --- Droite numérique]
Sur une droite graduée, chaque point a une abscisse, et chaque nombre est l'abscisse
d'un seul point. Les nombres ainsi repérés sont les \textbf{nombres réels} ; leur
ensemble est noté $\R$, et la droite graduée s'appelle la \textbf{droite numérique}.
\begin{center}
\begin{tikzpicture}[x=1.1cm]
  \axegrad{-3}{4}
  \foreach \p in {-2.5,-1,0,0.75,2,3.1416} {\path[coulDef] (\p,0) node[croix]{};}
  \node[above=2pt,font=\scriptsize] at (-2.5,0) {$-\frac52$};
  \node[above=2pt,font=\scriptsize] at (-1,0) {$-1$};
  \node[above=2pt,font=\scriptsize] at (0,0) {$0$};
  \node[above=2pt,font=\scriptsize] at (0.75,0) {$0{,}75$};
  \node[above=2pt,font=\scriptsize] at (2,0) {$2$};
  \node[above=2pt,font=\scriptsize] at (3.1416,0) {$\pi$};
\end{tikzpicture}
\end{center}
\end{definitionbox}

\begin{definitionbox}[Définition 2 --- Symboles de comparaison]
\begin{center}
\renewcommand{\arraystretch}{1.3}
\begin{tabular}{|c|l|l|}
\hline
\rowcolor{coulDef!10}\textbf{On écrit} & \textbf{On lit} & \textbf{Cela veut dire}\\\hline
$a<b$ & $a$ est \textbf{strictement inférieur} à $b$ & $a$ est plus petit que $b$, sans lui être égal\\\hline
$a\le b$ & $a$ est \textbf{inférieur ou égal} à $b$ & $a<b$ \textbf{ou} $a=b$\\\hline
$a>b$ & $a$ est \textbf{strictement supérieur} à $b$ & $a$ est plus grand que $b$, sans lui être égal\\\hline
$a\ge b$ & $a$ est \textbf{supérieur ou égal} à $b$ & $a>b$ \textbf{ou} $a=b$\\\hline
\end{tabular}
\end{center}
Le petit trait sous $\le$ et $\ge$ dit que l'égalité est permise : $4\le 4$ est vrai,
$4<4$ est faux ; $-1\ge -6$ est vrai.
\end{definitionbox}

\begin{exemplebox}[Exemple --- l'intervalle $\intff{-1}{3}$]
Les nombres $x$ compris entre $-1$ et $3$, ces deux nombres compris, vérifient
$-1\le x\le 3$. Leur ensemble est l'\textbf{intervalle} $\intff{-1}{3}$ ; $-1$ et $3$
en sont les \textbf{bornes}.
\begin{center}
\begin{tikzpicture}[x=0.8cm]
  \axegrad{-2}{5}
  \bornelab{-1}{-1}\bornelab{0}{0}\bornelab{1}{1}\bornelab{3}{3}
  \morceau{-1}{3}\crochet{-1}{[}\crochet{3}{]}
\end{tikzpicture}
\end{center}
On écrit $0\in\intff{-1}{3}$, $3\in\intff{-1}{3}$ et $4\notin\intff{-1}{3}$ :
$\in$ se lit \og appartient à\fg{}, $\notin$ \og n'appartient pas à\fg{}.
\end{exemplebox}

\begin{definitionbox}[Définition 3 --- Intervalle fermé]
Pour deux réels $a<b$, l'intervalle $\intff{a}{b}$ est l'ensemble des réels $x$ tels
que $a\le x\le b$. Ses deux bornes lui appartiennent : les crochets sont
\textbf{tournés vers l'intervalle}. On dit que l'intervalle est \textbf{fermé}.
\end{definitionbox}

\begin{methodebox}[Méthode 1 --- Représenter un intervalle fermé]
\etapeexemples
Représenter sur une droite graduée :
\begin{questions}[label=\alph*)]
  \item $I=\intff{-2}{3}$
  \item $J=\intff{-\frac32}{\frac12}$
\end{questions}
\begin{correction}
\cas{a)}\par\nopagebreak
\begin{center}
\begin{tikzpicture}[x=0.7cm]
  \axegrad{-3}{4}
  \bornelab{-2}{-2}\bornelab{0}{0}\bornelab{3}{3}
  \morceau{-2}{3}\crochet{-2}{[}\crochet{3}{]}
  \node[coulInt,right,font=\scriptsize] at (4.6,0) {$I$};
\end{tikzpicture}
\end{center}
\explique{On place les bornes $-2$ et $3$, on trace les crochets tournés vers
l'intérieur, puis on épaissit le trait entre les deux.}{les bornes sont dans $I$}
\cas{b)}\par\nopagebreak
\begin{center}
\begin{tikzpicture}[x=0.7cm]
  \axegrad{-3}{2}
  \bornelab{-1.5}{-\frac32}\bornelab{0}{0}\bornelab{0.5}{\frac12}
  \morceau{-1.5}{0.5}\crochet{-1.5}{[}\crochet{0.5}{]}
  \node[coulInt,right,font=\scriptsize] at (2.6,0) {$J$};
\end{tikzpicture}
\end{center}
\explique{$-\frac32=-1{,}5$ est au milieu de $-2$ et $-1$ ; $\frac12=0{,}5$ au milieu de
$0$ et $1$.}{une fraction se place par son écriture décimale}
\end{correction}

\etape{À vous de jouer}
Représenter sur une droite graduée : \quad $\intff{1}{6}$ \qquad $\intff{-5}{-2}$
\qquad $\intff{-\frac12}{\frac52}$
\end{methodebox}

\begin{methodebox}[Méthode 2 --- Dire si un nombre appartient à un intervalle]
\etapeexemples
Pour chaque nombre $1$ ; $\frac23$ ; $\frac12$ ; $\pi$ ; $4{,}01$, dire s'il appartient à
$I=\intff{\frac23}{4}$.
\begin{correction}
\cas{Résolution}\par\nopagebreak
\explique{$1\in I$, car $\frac23\le 1\le 4$.}{on compare aux deux bornes}
\explique{$\frac23\in I$ : c'est une borne, et le crochet est fermé.}{crochet tourné vers
l'intervalle}
\explique{$\frac12\notin I$, car $\frac12=0{,}5$ et $\frac23\approx 0{,}67$.}{on passe
aux écritures décimales}
\explique{$\pi\in I$, car $\pi\approx 3{,}14$.}{}
\explique{$4{,}01\notin I$, car $4{,}01>4$.}{}
\end{correction}

\etape{À vous de jouer}
\begin{questions}
  \item Avec $J=\intff{-3}{\frac52}$, compléter par $\in$ ou $\notin$ : \quad
        $-3\;\ldots\;J$ \qquad $\frac52\;\ldots\;J$ \qquad $2{,}6\;\ldots\;J$ \qquad
        $\pi\;\ldots\;J$ \qquad $-3{,}1\;\ldots\;J$
  \item Même question avec $K=\intff{-1}{\pi}$ : \quad $-1\;\ldots\;K$ \qquad
        $3\;\ldots\;K$ \qquad $3{,}2\;\ldots\;K$ \qquad $-1{,}5\;\ldots\;K$
\end{questions}
\reponse{1\textdegree) $-3\in J$ ;\; $\frac52\in J$ ;\; $2{,}6\notin J$ ;\; $\pi\notin J$ ;\;
$-3{,}1\notin J$.\; 2\textdegree) $-1\in K$ ;\; $3\in K$ ;\; $3{,}2\notin K$ ($\pi\approx
3{,}14$) ;\; $-1{,}5\notin K$.}
\end{methodebox}

\afaire{exercices 1 à 3}

\partiecours{Intervalles ouverts, inégalités}

\begin{definitionbox}[Définition 4 --- Bornes exclues, infini]
\begin{itemize}
  \item Un crochet \textbf{tourné vers l'extérieur} indique que la borne \textbf{n'appartient
        pas} à l'intervalle.
  \item Un intervalle sans aucune de ses bornes est \textbf{ouvert} ; avec une seule de
        ses bornes, il est \textbf{semi-ouvert}.
  \item Quand l'intervalle n'est pas limité d'un côté, on écrit $+\infty$ (\og plus
        l'infini\fg{}) ou $-\infty$ (\og moins l'infini\fg{}) ; de ce côté, le crochet
        est toujours tourné vers l'extérieur. L'infini n'est pas un nombre : on ne
        peut jamais l'atteindre.
\end{itemize}
\end{definitionbox}

\besoinplace{18\baselineskip}
\begin{exemplebox}[Exemples]
\renewcommand{\arraystretch}{1.6}
\begin{tabular}{@{}l@{\qquad}c@{\qquad}l@{}}
\textbf{Inégalité(s)} & \textbf{Représentation} & \textbf{Intervalle}\\
$-2<x<3$ &
\begin{tikzpicture}[x=0.45cm,baseline=-3pt]\axegrad{-4}{6}
\morceau{-2}{3}\crochet{-2}{]}\crochet{3}{[}\bornelab{-2}{-2}\bornelab{0}{0}\bornelab{3}{3}\end{tikzpicture}
& $x\in\intoo{-2}{3}$ \quad (ouvert)\\
$-2<x\le 3$ &
\begin{tikzpicture}[x=0.45cm,baseline=-3pt]\axegrad{-4}{6}
\morceau{-2}{3}\crochet{-2}{]}\crochet{3}{]}\bornelab{-2}{-2}\bornelab{0}{0}\bornelab{3}{3}\end{tikzpicture}
& $x\in\intof{-2}{3}$ \quad ($-2$ exclu, $3$ inclus)\\
$x>1$ &
\begin{tikzpicture}[x=0.45cm,baseline=-3pt]\axegrad{-4}{6}
\morceau{1}{6.5}\crochet{1}{]}\bornelab{0}{0}\bornelab{1}{1}\end{tikzpicture}
& $x\in\intoo{1}{+\infty}$\\
$x\le 4$ &
\begin{tikzpicture}[x=0.45cm,baseline=-3pt]\axegrad{-4}{6}
\morceau{-4.5}{4}\crochet{4}{]}\bornelab{0}{0}\bornelab{4}{4}\end{tikzpicture}
& $x\in\intof{-\infty}{4}$\\
tous les réels &
\begin{tikzpicture}[x=0.45cm,baseline=-3pt]\axegrad{-4}{6}
\morceau{-4.5}{6.5}\bornelab{0}{0}\end{tikzpicture}
& $x\in\intoo{-\infty}{+\infty}$, c'est-à-dire $x\in\R$
\end{tabular}
\end{exemplebox}

\begin{proprietebox}[Propriété 1 --- À retenir]
\begin{center}
\renewcommand{\arraystretch}{1.35}
\begin{tabular}{|>{\centering\arraybackslash}m{3.2cm}|>{\centering\arraybackslash}m{3.2cm}|>{\centering\arraybackslash}m{4.4cm}|}
\hline
\rowcolor{coulProp!10}\textbf{Réels $x$ tels que} & \textbf{Notation} & \textbf{Représentation}\\\hline
$a\le x<b$ & $x\in\intfo{a}{b}$ & \mini{\morceau{0}{3}\crochet{0}{[}\crochet{3}{[}\bornelab{0}{a}\bornelab{3}{b}}\\\hline
$x>a$ & $x\in\intoo{a}{+\infty}$ & \mini{\morceau{1}{5.6}\crochet{1}{]}\bornelab{1}{a}}\\\hline
$x\le b$ & $x\in\intof{-\infty}{b}$ & \mini{\morceau{-2.6}{3}\crochet{3}{]}\bornelab{3}{b}}\\\hline
\end{tabular}
\end{center}
$\le$ ou $\ge$ : crochet \textbf{fermé} ; \quad $<$, $>$, $+\infty$ ou $-\infty$ : crochet
\textbf{ouvert}.
\end{proprietebox}

\begin{methodebox}[Méthode 3 --- Passer d'une inégalité à un intervalle, et inversement]
\etapeexemples
\begin{questions}[label=\alph*)]
  \item Écrire avec un intervalle, puis représenter : \quad ${-1<x\le 4}$ \qquad
        ${x\ge -2}$ \qquad ${x<\frac32}$
  \item Écrire avec une inégalité : \quad ${x\in\intfo{-5}{1}}$ \qquad
        ${x\in\intoo{2}{+\infty}}$
\end{questions}
\begin{correction}
\cas{a)}\par\nopagebreak
\begin{tabular}[t]{@{}l@{\qquad}l@{}}
$-1<x\le 4$ \;:\; $x\in\intof{-1}{4}$ &
\begin{tikzpicture}[x=0.45cm,baseline=-3pt]\axegrad{-4}{6}
\morceau{-1}{4}\crochet{-1}{]}\crochet{4}{]}\bornelab{-1}{-1}\bornelab{0}{0}\bornelab{4}{4}\end{tikzpicture}\\[12pt]
$x\ge -2$ \;:\; $x\in\intfo{-2}{+\infty}$ &
\begin{tikzpicture}[x=0.45cm,baseline=-3pt]\axegrad{-4}{6}
\morceau{-2}{6.5}\crochet{-2}{[}\bornelab{-2}{-2}\bornelab{0}{0}\end{tikzpicture}\\[12pt]
$x<\frac32$ \;:\; $x\in\intoo{-\infty}{\frac32}$ &
\begin{tikzpicture}[x=0.45cm,baseline=-3pt]\axegrad{-4}{6}
\morceau{-4.5}{1.5}\crochet{1.5}{[}\bornelab{0}{0}\bornelab{1.5}{\frac32}\end{tikzpicture}
\end{tabular}
\explique{Les crochets suivent les symboles : fermé pour $\le$ et $\ge$, ouvert pour $<$,
$>$ et du côté de l'infini.}{on écrit toujours la plus petite borne à gauche}
\cas{b)}
\explique{$x\in\intfo{-5}{1}$ s'écrit $-5\le x<1$.}{$-5$ inclus, $1$ exclu}
\explique{$x\in\intoo{2}{+\infty}$ s'écrit $x>2$.}{$2$ exclu}
\end{correction}

\etape{À vous de jouer}
\begin{questions}
  \item Écrire avec un intervalle, puis représenter : \quad $2\le x<7$ \qquad $x>-3$
        \qquad $x\le 1{,}5$ \qquad $-2<x<\frac12$
  \item Écrire avec une inégalité : \quad $x\in\intoo{-\infty}{4}$ \qquad
        $x\in\intof{-3}{2}$ \qquad $x\in\intfo{0}{+\infty}$
\end{questions}
\reponse{1\textdegree) $\intfo{2}{7}$ ;\; $\intoo{-3}{+\infty}$ ;\; $\intof{-\infty}{1{,}5}$ ;\;
$\intoo{-2}{\frac12}$.\; 2\textdegree) $x<4$ ;\; $-3<x\le 2$ ;\; $x\ge 0$.}
\end{methodebox}

\afaire{exercices 4 à 8}

\partiecours{Intersection, réunion, complémentaire}

\begin{definitionbox}[Définition 5 --- Intersection]
L'\textbf{intersection} de deux intervalles $I$ et $J$, notée $I\cap J$ (\og $I$ inter
$J$\fg{}), est l'ensemble des nombres qui appartiennent \textbf{à la fois} à $I$ et à
$J$. S'il n'y en a aucun, l'intersection est l'\textbf{ensemble vide}, noté
$\varnothing$.
\end{definitionbox}

\begin{exemplebox}[Exemple --- $I=\intff{-1}{3}$ et $J=\intff{1}{5}$]
\begin{center}
\begin{tikzpicture}[x=0.8cm]
  \surligne{1}{3}
  \axegrad{-2}{6}
  \foreach \k in {-1,0,1,3,5} {\bornelab[-0.3]{\k}{\k}}
  \morceau[coulInt][0.3]{-1}{3}\crochet[coulInt][0.3]{-1}{[}\crochet[coulInt][0.3]{3}{]}
  \morceau[coulIntB][-0.3]{1}{5}\crochet[coulIntB][-0.3]{1}{[}\crochet[coulIntB][-0.3]{5}{]}
  \node[coulInt,right,font=\scriptsize] at (6.6,0.3) {$I$};
  \node[coulIntB,right,font=\scriptsize] at (6.6,-0.3) {$J$};
\end{tikzpicture}
\end{center}
Les deux intervalles se superposent de $1$ à $3$ (zone surlignée), et $1$ et $3$ sont
dans $I$ et dans $J$ : \quad $I\cap J=\intff{1}{3}$.
\end{exemplebox}

\begin{definitionbox}[Définition 6 --- Réunion]
La \textbf{réunion} de deux intervalles $I$ et $J$, notée $I\cup J$ (\og $I$ union
$J$\fg{}), est l'ensemble des nombres qui appartiennent à $I$ \textbf{ou} à $J$ :
à $I$ seulement, à $J$ seulement, ou aux deux.
\end{definitionbox}

\begin{exemplebox}[Exemple --- $I=\intff{-1}{3}$ et $J=\intff{1}{5}$]
\begin{center}
\begin{tikzpicture}[x=0.8cm]
  \surligne{-1}{5}
  \axegrad{-2}{6}
  \foreach \k in {-1,0,1,3,5} {\bornelab[-0.3]{\k}{\k}}
  \morceau[coulInt][0.3]{-1}{3}\crochet[coulInt][0.3]{-1}{[}\crochet[coulInt][0.3]{3}{]}
  \morceau[coulIntB][-0.3]{1}{5}\crochet[coulIntB][-0.3]{1}{[}\crochet[coulIntB][-0.3]{5}{]}
  \node[coulInt,right,font=\scriptsize] at (6.6,0.3) {$I$};
  \node[coulIntB,right,font=\scriptsize] at (6.6,-0.3) {$J$};
\end{tikzpicture}
\end{center}
À eux deux, les intervalles couvrent la droite de $-1$ à $5$ (zone surlignée) :
\quad $I\cup J=\intff{-1}{5}$.
\end{exemplebox}

\begin{methodebox}[Méthode 4 --- Trouver l'intersection et la réunion]
\etapeexemples
Déterminer $I\cap J$ et $I\cup J$ :
\begin{questions}[label=\alph*)]
  \item $I=\intff{-2}{2}$ et $J=\intoo{0}{5}$
  \item $I=\intof{-\infty}{-2}$ et $J=\intff{1}{3}$
\end{questions}
\begin{correction}
\cas{a)}\par\nopagebreak
\begin{center}\begin{tikzpicture}[x=0.6cm]
  \axegrad{-3}{6}
  \bornelab[-0.3]{-2}{-2}\bornelab[-0.3]{0}{0}\bornelab[-0.3]{2}{2}\bornelab[-0.3]{5}{5}
  \morceau[coulInt][0.3]{-2}{2}\crochet[coulInt][0.3]{-2}{[}\crochet[coulInt][0.3]{2}{]}
  \morceau[coulIntB][-0.3]{0}{5}\crochet[coulIntB][-0.3]{0}{]}\crochet[coulIntB][-0.3]{5}{[}
  \node[coulInt,right,font=\scriptsize] at (6.6,0.3) {$I$};
  \node[coulIntB,right,font=\scriptsize] at (6.6,-0.3) {$J$};
\end{tikzpicture}\end{center}
\explique{Les intervalles se superposent de $0$ à $2$ ; $0\notin J$ et
$2\in I\cap J$ : \resultat{I\cap J=\intof{0}{2}}.}{on représente $I$ et $J$ sur le même
axe}
\explique{Ils couvrent de $-2$ ($-2\in I$) à $5$ ($5$ n'est ni dans $I$ ni dans $J$) :
\resultat{I\cup J=\intfo{-2}{5}}.}{une borne de $I\cup J$ est dans $I$ ou dans $J$}
\cas{b)}\par\nopagebreak
\begin{center}\begin{tikzpicture}[x=0.6cm]
  \axegrad{-4}{5}
  \bornelab[-0.3]{-2}{-2}\bornelab[-0.3]{0}{0}\bornelab[-0.3]{1}{1}\bornelab[-0.3]{3}{3}
  \morceau[coulInt][0.3]{-4.5}{-2}\crochet[coulInt][0.3]{-2}{]}
  \morceau[coulIntB][-0.3]{1}{3}\crochet[coulIntB][-0.3]{1}{[}\crochet[coulIntB][-0.3]{3}{]}
  \node[coulInt,right,font=\scriptsize] at (5.6,0.3) {$I$};
  \node[coulIntB,right,font=\scriptsize] at (5.6,-0.3) {$J$};
\end{tikzpicture}\end{center}
\explique{Aucune zone commune : \resultat{I\cap J=\varnothing}.}{}
\explique{La réunion ne s'écrit pas avec un seul intervalle :
\resultat{I\cup J=\intof{-\infty}{-2}\cup\intff{1}{3}}.}{il y a un trou entre $-2$ et
$1$}
\end{correction}

\etape{À vous de jouer}
Déterminer $I\cap J$ et $I\cup J$ :
\begin{questions}[label=\alph*)]
  \item $I=\intfo{-3}{2}$ et $J=\intfo{-1}{+\infty}$
  \item $I=\intof{-\infty}{1}$ et $J=\intof{1}{4}$
  \item $I=\intff{0}{5}$ et $J=\intoo{1}{3}$
\end{questions}
\reponse{a) $I\cap J=\intfo{-1}{2}$, $I\cup J=\intfo{-3}{+\infty}$ ;\; b) $I\cap
J=\varnothing$ ($1\notin J$), $I\cup J=\intof{-\infty}{4}$ ;\; c) $I\cap J=\intoo{1}{3}$,
$I\cup J=\intff{0}{5}$ ($J$ est dans $I$).}
\end{methodebox}

\begin{definitionbox}[Définition 7 --- Complémentaire]
Le \textbf{complémentaire} dans $\R$ d'un intervalle $I$, noté $\complement_{\R} I$, est
l'ensemble des réels qui \textbf{n'appartiennent pas} à $I$.
\end{definitionbox}

\begin{exemplebox}[Exemple --- $I=\intfo{2}{4}$]
\begin{center}
\begin{tikzpicture}[x=0.8cm]
  \surligne{-1.5}{2}\surligne{4}{7.5}
  \axegrad{-1}{7}
  \foreach \k in {0,1,2,4} {\bornelab[-0.3]{\k}{\k}}
  \morceau[coulInt][0.3]{2}{4}\crochet[coulInt][0.3]{2}{[}\crochet[coulInt][0.3]{4}{[}
  \morceau[coulIntB][-0.3]{-1.5}{2}\crochet[coulIntB][-0.3]{2}{[}
  \morceau[coulIntB][-0.3]{4}{7.5}\crochet[coulIntB][-0.3]{4}{[}
  \node[coulInt,right,font=\scriptsize] at (7.6,0.3) {$I$};
  \node[coulIntB,right,font=\scriptsize] at (7.6,-0.3) {$\complement_{\R} I$};
\end{tikzpicture}
\end{center}
Le complémentaire est tout le reste de la droite. $2\in I$, donc $2$ n'est pas dans le
complémentaire ; $4\notin I$, donc $4$ y est : \quad
$\complement_{\R} I=\intoo{-\infty}{2}\cup\intfo{4}{+\infty}$.
\end{exemplebox}

\begin{methodebox}[Méthode 5 --- Trouver le complémentaire d'un intervalle]
\etapeexemples
Déterminer le complémentaire dans $\R$ de :
\begin{questions}[label=\alph*)]
  \item $I=\intff{1}{3}$
  \item $K=\intoo{-\infty}{2}$
\end{questions}
\begin{correction}
\cas{a)}\par\nopagebreak
\begin{center}\begin{tikzpicture}[x=0.6cm]
  \axegrad{-2}{6}
  \bornelab[-0.3]{0}{0}\bornelab[-0.3]{1}{1}\bornelab[-0.3]{3}{3}
  \morceau[coulInt][0.3]{1}{3}\crochet[coulInt][0.3]{1}{[}\crochet[coulInt][0.3]{3}{]}
  \morceau[coulIntB][-0.3]{-2.5}{1}\crochet[coulIntB][-0.3]{1}{[}
  \morceau[coulIntB][-0.3]{3}{6.5}\crochet[coulIntB][-0.3]{3}{]}
  \node[coulInt,right,font=\scriptsize] at (6.6,0.3) {$I$};
  \node[coulIntB,right,font=\scriptsize] at (6.6,-0.3) {$\complement_{\R} I$};
\end{tikzpicture}\end{center}
\explique{$1$ et $3$ sont dans $I$, donc pas dans son complémentaire :
\resultat{\complement_{\R} I=\intoo{-\infty}{1}\cup\intoo{3}{+\infty}}.}{chaque crochet
se retourne}
\cas{b)}
\explique{$2\notin K$, donc $2$ est dans le complémentaire :
\resultat{\complement_{\R} K=\intfo{2}{+\infty}}.}{le crochet en $2$ se retourne}
\end{correction}

\etape{À vous de jouer}
Déterminer le complémentaire dans $\R$ de : \quad $L=\intoo{-2}{+\infty}$ \qquad
$M=\intfo{1}{4}$ \qquad $N=\intof{-\infty}{3}$
\reponse{$\complement_{\R} L=\intof{-\infty}{-2}$ ;\;
$\complement_{\R} M=\intoo{-\infty}{1}\cup\intfo{4}{+\infty}$ ;\;
$\complement_{\R} N=\intoo{3}{+\infty}$.}
\end{methodebox}

\afaire{exercices 9 à 13}

\partiecours{Image et antécédent}

\begin{definitionbox}[Définition 8 --- Fonction, image, antécédent]
Une \textbf{fonction} $f$ associe à chaque nombre $x$ de départ \textbf{un seul} nombre,
noté $f(x)$ (on lit \og $f$ de $x$\fg{}).
\begin{itemize}
  \item $f(x)$ est l'\textbf{image} de $x$ par $f$ ;
  \item $x$ est \textbf{un antécédent} de $f(x)$ par $f$.
\end{itemize}
\begin{center}
\begin{tikzpicture}[>={Stealth[length=2.5mm]}]
  \node (a) at (0,0) {$x$};
  \node[draw=coulDef,thick,fill=coulDef!8,minimum width=1.4cm,minimum height=0.75cm,rounded corners] (f) at (2.4,0) {$f$};
  \node (b) at (4.8,0) {$f(x)$};
  \draw[->,thick] (a)--(f); \draw[->,thick] (f)--(b);
  \node[below,font=\scriptsize,coulDef] at (a.south) {antécédent};
  \node[below,font=\scriptsize,coulDef] at (b.south) {image};
\end{tikzpicture}
\end{center}
Exemple : avec $f(x)=3x-2$, on a $f(4)=3\times 4-2=10$ ; $10$ est l'image de $4$, et $4$
est un antécédent de $10$. On écrit aussi $f : 4\longmapsto 10$.
\end{definitionbox}

\begin{methodebox}[Méthode 6 --- Calculer une image]
On remplace $x$ par le nombre, \textbf{entre parenthèses} s'il est négatif, puis on
calcule en respectant les priorités.
\etapeexemples
\begin{questions}[label=\alph*)]
  \item Soit ${f(x)=4x-3}$. Calculer $f(3)$ et $f(-2)$.
  \item Soit ${h(x)=2x^{2}-3x+1}$. Calculer $h(3)$, $h(-2)$ et $h\!\left(\frac12\right)$.
\end{questions}
\begin{correction}
\cas{a)}
\begin{calculs}
f(3) &= 4\times 3-3 &\qquad& \rem{on remplace $x$ par $3$}\\
f(3) &= 9 &&\\
f(-2) &= 4\times(-2)-3 &\qquad& \rem{$-2$ entre parenthèses}\\
f(-2) &= -11 &&
\end{calculs}
\cas{b)}
\begin{calculs}
h(3) &= 2\times 3^{2}-3\times 3+1 &\qquad& \rem{on remplace $x$ par $3$}\\
h(3) &= 18-9+1 &\qquad& \rem{la puissance d'abord : $3^{2}=9$}\\
h(3) &= 10 &&\\
h(-2) &= 2\times(-2)^{2}-3\times(-2)+1 &\qquad& \rem{$(-2)^{2}=4$, et non $-4$}\\
h(-2) &= 8+6+1 &&\\
h(-2) &= 15 &&\\
h\!\left(\tfrac12\right) &= 2\times\left(\tfrac12\right)^{2}-3\times\tfrac12+1 &\qquad& \rem{$\left(\frac12\right)^{2}=\frac14$}\\
h\!\left(\tfrac12\right) &= \tfrac12-\tfrac32+1 &&\\
h\!\left(\tfrac12\right) &= 0 &\qquad& \rem{$\frac12-\frac32=-1$}
\end{calculs}
\end{correction}

\etape{À vous de jouer}
\begin{questions}
  \item Soit ${g(x)=\frac23x+1}$. Calculer $g(6)$, $g(0)$ et $g\!\left(\frac34\right)$.
  \item Soit ${k(x)=-x^{2}+3x+2}$. Calculer $k(2)$, $k(-1)$ et $k\!\left(\frac12\right)$.
\end{questions}
\reponse{1\textdegree) $g(6)=5$ ;\; $g(0)=1$ ;\; $g\!\left(\frac34\right)=\frac12+1=\frac32$.\;
2\textdegree) $k(2)=-4+6+2=4$ ;\; $k(-1)=-1-3+2=-2$ ;\;
$k\!\left(\frac12\right)=-\frac14+\frac32+2=\frac{13}{4}$.}
\end{methodebox}

\begin{methodebox}[Méthode 7 --- Chercher ou vérifier un antécédent]
Chercher les antécédents de $b$, c'est \textbf{résoudre l'équation} $f(x)=b$. Vérifier
que $a$ est un antécédent de $b$, c'est \textbf{calculer} $f(a)$ et regarder si l'on
trouve $b$.
\etapeexemples
\begin{questions}[label=\alph*)]
  \item Soit ${f(x)=3x+2}$. Déterminer l'antécédent de $-10$ par $f$.
  \item Soit ${p(x)=x^{2}+2x}$. Le nombre $-3$ est-il un antécédent de $3$ par $p$ ?
        Et le nombre $2$ ?
\end{questions}
\begin{correction}
\cas{a)}
\begin{calculs}
3x+2 &= -10 &\qquad& \rem{on résout $f(x)=-10$}\\
3x &= -12 &\qquad& \rem{on retranche $2$}\\
x &= -4 &\qquad& \rem{on divise par $3$}
\end{calculs}
\explique{L'antécédent de $-10$ par $f$ est $-4$.}{vérification : $3\times(-4)+2=-10$}
\cas{b)}
\begin{calculs}
p(-3) &= (-3)^{2}+2\times(-3) &\qquad& \rem{on calcule l'image de $-3$}\\
p(-3) &= 9-6=3 &&\\
p(2) &= 2^{2}+2\times 2=8 &&
\end{calculs}
\explique{$p(-3)=3$ : $-3$ est un antécédent de $3$ ; $p(2)=8\ne 3$ : $2$ n'en est
pas un.}{on compare au nombre $3$}
\end{correction}

\etape{À vous de jouer}
\begin{questions}
  \item Soit ${u(x)=4x-7}$. Déterminer l'antécédent de $9$ par $u$.
  \item Soit ${k(x)=-x^{2}+3x+2}$. Le nombre $-1$ est-il un antécédent de $-2$ par $k$ ?
        Et le nombre $3$ ?
\end{questions}
\reponse{1\textdegree) $4x-7=9$, donc $x=4$.\; 2\textdegree) $k(-1)=-2$ : oui ;\;
$k(3)=-9+9+2=2\ne -2$ : non.}
\end{methodebox}

\afaire{exercices 14 à 17}

\partiecours{Ensemble de définition}

\begin{exemplebox}[Deux calculs impossibles]
Certains nombres n'ont pas d'image, parce que le calcul est impossible : ce sont des
\textbf{valeurs interdites}.
\begin{itemize}
  \item Avec $f(x)=\dfrac{1}{x}$, $f(0)=1\div 0$ n'existe pas : \textbf{on ne divise pas
        par zéro}. Le nombre $0$ est une valeur interdite.
  \item Avec $g(x)=\sqrt{x}$, $g(-4)=\sqrt{-4}$ n'existe pas : \textbf{un nombre négatif
        n'a pas de racine carrée}. Tous les nombres strictement négatifs sont interdits.
\end{itemize}
\end{exemplebox}

\begin{definitionbox}[Définition 9 --- Ensemble de définition]
L'\textbf{ensemble de définition} d'une fonction $f$, noté $D_f$, est l'ensemble des
nombres qui ont une image par $f$. On note $f : D_f\longrightarrow\R$,
$x\longmapsto f(x)$.

\smallskip
Exemples : $D_f=\intoo{-\infty}{0}\cup\intoo{0}{+\infty}$ pour $f(x)=\frac1x$ ;
$D_g=\intfo{0}{+\infty}$ pour $g(x)=\sqrt{x}$.
\end{definitionbox}

\begin{proprietebox}[Propriété 2 --- À retenir]
\begin{itemize}
  \item Une \textbf{fonction affine} $f(x)=ax+b$ est définie sur $\R$, comme toute
        expression sans quotient ni racine carrée (par exemple $x^{2}+5x-3$).
  \item Un \textbf{dénominateur} ne doit pas être nul : les nombres qui l'annulent sont
        des valeurs interdites.
  \item Sous une \textbf{racine carrée}, le nombre doit être \textbf{positif ou nul}.
\end{itemize}
\end{proprietebox}

\begin{methodebox}[Méthode 8 --- Trouver la valeur interdite d'un quotient]
On résout l'équation \og dénominateur $=0$\fg{}.
\etapeexemples
Trouver la valeur interdite de \quad $h(x)=\dfrac{1}{3x-2}$ \quad et de \quad
$u(x)=\dfrac{x+1}{x-4}$.
\begin{correction}
\cas{Résolution}\par\nopagebreak
\begin{calculs}
3x-2 &= 0 &\qquad& \rem{le dénominateur de $h$ s'annule}\\
3x &= 2 &&\\
x &= \tfrac23 &&\\
x-4 &= 0 &\qquad& \rem{le dénominateur de $u$ s'annule}\\
x &= 4 &&
\end{calculs}
\explique{La valeur interdite de $h$ est $\frac23$, celle de $u$ est $4$.}{$h\!\left(\frac23\right)$
et $u(4)$ n'existent pas}
\end{correction}

\etape{À vous de jouer}
Trouver la valeur interdite de \quad $p(x)=\dfrac{x-2}{2x+8}$ \quad et de \quad
$v(x)=\dfrac{3}{x+5}$.
\reponse{$2x+8=0$ : valeur interdite $-4$ ;\; $x+5=0$ : valeur interdite $-5$.}
\end{methodebox}

\begin{methodebox}[Méthode 9 --- Déterminer un ensemble de définition]
Sans quotient ni racine : $D_f=\R$. Avec un quotient : on retire de $\R$ la valeur
interdite. Avec une racine : on résout \og nombre sous la racine $\ge 0$\fg{} comme une
équation (on ajoute ou on retranche un même nombre, on divise par un même nombre
positif).
\etapeexemples
Déterminer l'ensemble de définition de :
\[
  f(x)=2x+7 \qquad h(x)=\dfrac{1}{3x-2} \qquad i(x)=\sqrt{x-4} \qquad j(x)=\sqrt{3x+6}
\]
\begin{correction}
\cas{Résolution}\par\nopagebreak
\explique{$f$ est affine : $D_f=\R$.}{ni quotient ni racine}
\explique{La valeur interdite de $h$ est $\frac23$ (méthode 8) :
$D_h=\intoo{-\infty}{\frac23}\cup\intoo{\frac23}{+\infty}$.}{on retire $\frac23$}
\begin{calculs}
x-4 &\ge 0 &\qquad& \rem{le nombre sous la racine de $i$ est positif ou nul}\\
x &\ge 4 &\qquad& \rem{on ajoute $4$}\\
3x+6 &\ge 0 &\qquad& \rem{même travail pour $j$}\\
3x &\ge -6 &\qquad& \rem{on retranche $6$}\\
x &\ge -2 &\qquad& \rem{on divise par $3$, qui est positif}
\end{calculs}
\explique{$D_i=\intfo{4}{+\infty}$ et $D_j=\intfo{-2}{+\infty}$.}{on conclut avec des
intervalles}
\end{correction}

\etape{À vous de jouer}
Déterminer l'ensemble de définition de :
\[
  g(x)=-2x+9 \qquad p(x)=\dfrac{x-2}{2x+8} \qquad q(x)=\sqrt{x+3} \qquad
  r(x)=\sqrt{2x-10} \qquad s(x)=4x-1
\]
\reponse{$D_g=\R$ ;\; $D_p=\intoo{-\infty}{-4}\cup\intoo{-4}{+\infty}$ ;\;
$D_q=\intfo{-3}{+\infty}$ ;\; $D_r=\intfo{5}{+\infty}$ ;\; $D_s=\R$.}
\end{methodebox}

\afaire{exercices 18 à 21}

\partiecours{Tableau de valeurs}

\begin{definitionbox}[Définition 10 --- Tableau de valeurs]
Un \textbf{tableau de valeurs} d'une fonction $f$ donne sur la première ligne des valeurs
de $x$ (des \textbf{antécédents}) et sur la seconde leurs \textbf{images} $f(x)$. Une
image se lit de haut en bas, un antécédent de bas en haut.

\smallskip
Le \textbf{pas} est l'écart entre deux valeurs de $x$ qui se suivent. De $-2$ à $2$ avec un
pas de $1$ : $-2$ ; $-1$ ; $0$ ; $1$ ; $2$. Avec un pas de $0{,}5$ : $-2$ ; $-1{,}5$ ;
$-1$ ; \ldots ; $2$.
\end{definitionbox}

\begin{methodebox}[Méthode 10 --- Dresser et lire un tableau de valeurs]
\etapeexemples
\begin{questions}[label=\alph*)]
  \item Soit ${f(x)=3x-2}$. Dresser le tableau de valeurs de $f$ de $-3$ à $3$ avec un pas
        de $1$. Donner l'image de $-3$ et un antécédent de $4$.
  \item Même travail avec $g(x)=x^{2}-1$, de $-3$ à $3$ avec un pas de $1$. Combien le
        nombre $3$ a-t-il d'antécédents dans le tableau ?
\end{questions}
\begin{correction}
\cas{a)}
\explique{$f(-3)=3\times(-3)-2=-11$, $f(-2)=-8$, et ainsi de suite.}{on calcule chaque
image (méthode 6)}
\begin{center}
\renewcommand{\arraystretch}{1.3}
\begin{tabular}{|c|*{7}{>{\centering\arraybackslash}p{0.75cm}|}l}
\cline{1-8}
$x$ & $-3$ & $-2$ & $-1$ & $0$ & $1$ & $2$ & $3$ & $\leftarrow$ antécédents\\\cline{1-8}
$f(x)$ & $-11$ & $-8$ & $-5$ & $-2$ & $1$ & $4$ & $7$ & $\leftarrow$ images\\\cline{1-8}
\end{tabular}
\end{center}
\explique{L'image de $-3$ est $-11$ ; un antécédent de $4$ est $2$.}{haut vers bas, puis
bas vers haut}
\cas{b)}\par\nopagebreak
\begin{center}
\renewcommand{\arraystretch}{1.3}
\begin{tabular}{|c|*{7}{>{\centering\arraybackslash}p{0.75cm}|}}
\hline
$x$ & $-3$ & $-2$ & $-1$ & $0$ & $1$ & $2$ & $3$\\\hline
$g(x)$ & $8$ & $3$ & $0$ & $-1$ & $0$ & $3$ & $8$\\\hline
\end{tabular}
\end{center}
\explique{$g(-2)=3$ et $g(2)=3$ : le nombre $3$ a deux antécédents, $-2$ et $2$.}{un
nombre peut avoir plusieurs antécédents}
\end{correction}

\begin{remarquebox}[À retenir]
Un nombre de départ n'a \textbf{qu'une seule image} ; un nombre peut avoir
\textbf{plusieurs antécédents}, ou aucun.
\end{remarquebox}

\etape{À vous de jouer}
Soit ${u(x)=-x^{2}+4x}$.
\begin{questions}
  \item Dresser le tableau de valeurs de $u$ de $-1$ à $5$ avec un pas de $1$.
  \item Donner l'image de $2$, puis les antécédents de $0$ et de $3$ lus dans le tableau.
\end{questions}
\reponse{1\textdegree) $u(-1)=-5$ ;\; $u(0)=0$ ;\; $u(1)=3$ ;\; $u(2)=4$ ;\; $u(3)=3$ ;\;
$u(4)=0$ ;\; $u(5)=-5$.\; 2\textdegree) $u(2)=4$ ;\; antécédents de $0$ : $0$ et $4$ ;\;
de $3$ : $1$ et $3$.}
\end{methodebox}

\afaire{exercices 22 à 24}

\partiecours{Courbe représentative}

\begin{definitionbox}[Définition 11 --- Courbe d'une fonction]
Dans un repère, la \textbf{courbe représentative} $\Cf$ d'une fonction $f$ est formée des
points de coordonnées $\bigl(x\,;f(x)\bigr)$ : antécédents en abscisse, images en
ordonnée. Un point $M(x\,;y)$ est sur $\Cf$ lorsque $y=f(x)$.
\end{definitionbox}

\begin{methodebox}[Méthode 11 --- Tracer une courbe, lire images et antécédents]
On place les points d'un tableau de valeurs et on les relie à main levée par une courbe
\textbf{sans angle}. L'\textbf{image} de $a$ se lit en ordonnée, au-dessus (ou
au-dessous) de $a$ ; les \textbf{antécédents} de $b$ sont les abscisses des points de la
courbe situés sur la droite horizontale $y=b$.
\etapeexemples
Soit ${f(x)=x^{2}-4}$, de $-3$ à $3$.
\begin{questions}[label=\alph*)]
  \item Dresser le tableau de valeurs de $f$ de $-3$ à $3$ avec un pas de $1$, puis
        tracer la courbe de $f$.
  \item Lire les antécédents de $5$, puis ceux de $-5$.
\end{questions}
\begin{correction}
\begin{minipage}[t]{0.58\linewidth}\raggedright
\cas{a)}\par\nopagebreak
\begin{center}
\renewcommand{\arraystretch}{1.25}
\begin{tabular}{|c|*{7}{c|}}
\hline
$x$ & $-3$ & $-2$ & $-1$ & $0$ & $1$ & $2$ & $3$\\\hline
$f(x)$ & $5$ & $0$ & $-3$ & $-4$ & $-3$ & $0$ & $5$\\\hline
\end{tabular}
\end{center}
On place les sept points, puis on les relie : la courbe obtenue est une
\textbf{parabole}.

\cas{b)}\par\nopagebreak
La droite $y=5$ coupe la courbe aux points d'abscisses $-3$ et $3$ : les antécédents
de $5$ sont ${-3}$ et ${3}$. La droite $y=-5$ ne coupe pas la courbe : $-5$ n'a pas
d'antécédent.
\end{minipage}\hfill
\begin{minipage}[t]{0.38\linewidth}
\vspace{0pt}\centering
\begin{tikzpicture}[x=0.42cm,y=0.42cm]
  \repere{-4}{-5}{4}{6}
  \gradx{-3,-2,-1,1,2,3}
  \grady{-4,-3,-2,-1,1,2,3,4,5}
  \draw[coulCourbe,line width=1.1pt,domain=-3:3,samples=60] plot (\x,{(\x)*(\x)-4});
  \foreach \a/\b in {-3/5,-2/0,-1/-3,0/-4,1/-3,2/0,3/5} \path[coulCourbe] (\a,\b) node[croix]{};
  \node[coulCourbe,font=\footnotesize,right] at (3,4.4) {$\Cf$};
  \draw[black!70,thick] (-4,5)--(4,5);
  \node[black!70,font=\scriptsize,above] at (-2.6,5) {$y=5$};
  \draw[black!60,thick] (-4,-5)--(4,-5);
  \node[black!60,font=\scriptsize,above] at (2.6,-5) {$y=-5$};
\end{tikzpicture}
\end{minipage}
\end{correction}

\begin{remarquebox}
Une courbe représente une fonction seulement si chaque droite verticale la coupe
\textbf{au plus une fois} : un nombre n'a qu'une image.
\end{remarquebox}
\end{methodebox}

\begin{methodebox}[Méthode 12 --- Savoir si un point est sur une courbe]
On calcule l'image de l'abscisse du point et on la compare à son ordonnée :
$M(x_M\,;y_M)$ est sur $\Cf$ lorsque $f(x_M)=y_M$.
\etapeexemples
Avec $f(x)=x^{2}-4$ : les points $A(2{,}5\,;2{,}25)$ et $B(-1\,;3)$ sont-ils sur $\Cf$ ?
\begin{correction}
\cas{Point $A$}\par\nopagebreak
\begin{calculs}
f(2{,}5) &= 2{,}5^{2}-4 &\qquad& \rem{l'image de l'abscisse de $A$}\\
f(2{,}5) &= 6{,}25-4 &&\\
f(2{,}5) &= 2{,}25 &&
\end{calculs}
\explique{$f(2{,}5)=2{,}25$, l'ordonnée de $A$ : $A\in\Cf$.}{on compare à l'ordonnée}
\cas{Point $B$}\par\nopagebreak
\begin{calculs}
f(-1) &= (-1)^{2}-4 &\qquad& \rem{l'image de l'abscisse de $B$}\\
f(-1) &= 1-4 &&\\
f(-1) &= -3 &&
\end{calculs}
\explique{$f(-1)=-3$, alors que l'ordonnée de $B$ est $3$ : $B\notin\Cf$.}{$-3\ne 3$}
\end{correction}

\etape{À vous de jouer}
\begin{minipage}[t]{0.56\linewidth}\raggedright
\begin{questions}
  \item Ci-contre, la courbe d'une fonction $f$. Lire :
        \begin{questions}
          \item l'ensemble de définition de $f$ ;
          \item l'image de $-3$, de $0$ et de $4$ ;
          \item les antécédents de $1$, de $-1$, de $2$ et de $4$.
        \end{questions}
  \item Soit ${u(x)=2x^{2}-5}$. Les points $A(-2\,;3)$, $B(1{,}5\,;-0{,}5)$ et
        $C(0\,;5)$ sont-ils sur la courbe de $u$ ?
\end{questions}
\end{minipage}\hfill
\begin{minipage}[t]{0.40\linewidth}
\vspace{0pt}\centering
\begin{tikzpicture}[x=0.4cm,y=0.4cm]
  \repere{-6}{-3}{5}{4}
  \gradx{-5,-4,-3,-2,-1,1,2,3,4}
  \grady{-2,-1,1,2,3}
  \draw[coulCourbe,line width=1.1pt] (-5,-1)
    .. controls (-4.667,-0.333) and (-4.333,0.5) .. (-4,1)
    .. controls (-3.667,1.5) and (-3.333,2) .. (-3,2)
    .. controls (-2.667,2) and (-2.333,1.5) .. (-2,1)
    .. controls (-1.667,0.5) and (-1.333,-0.5) .. (-1,-1)
    .. controls (-0.667,-1.5) and (-0.333,-2) .. (0,-2)
    .. controls (0.333,-2) and (0.667,-1.5) .. (1,-1)
    .. controls (1.333,-0.5) and (1.667,0.5) .. (2,1)
    .. controls (2.333,1.5) and (2.667,1.667) .. (3,2)
    .. controls (3.333,2.333) and (3.667,2.667) .. (4,3);
  \path[coulCourbe] (-5,-1) node[croix]{} (4,3) node[croix]{};
  \node[coulCourbe,font=\footnotesize] at (-4.4,2.5) {$\Cf$};
\end{tikzpicture}
\end{minipage}
\reponse{1\textdegree) a) $D_f=\intff{-5}{4}$ ;\; b) $f(-3)=2$, $f(0)=-2$, $f(4)=3$ ;\;
c) antécédents de $1$ : $-4$, $-2$ et $2$ ; de $-1$ : $-5$, $-1$ et $1$ ; de $2$ : $-3$
et $3$ ; $4$ n'en a pas.\; 2\textdegree) $u(-2)=3$ : $A$ oui ;\; $u(1{,}5)=-0{,}5$ : $B$
oui ;\; $u(0)=-5\ne 5$ : $C$ non.}
\end{methodebox}

\afaire{exercices 25 à 28}

\partiecours{Calculatrice NumWorks}

\begin{proprietebox}[Application \texttt{Calculs}]
\begin{center}
\renewcommand{\arraystretch}{1.3}
\begin{tabular}{|>{\raggedright\arraybackslash}p{4.4cm}|>{\raggedright\arraybackslash}p{6.4cm}|>{\raggedright\arraybackslash}p{5.4cm}|}
\hline
\rowcolor{coulProp!10}\textbf{Pour} & \textbf{On tape} & \textbf{On obtient}\\\hline
une puissance & $3$ \fbox{\small$x^{y}$} $4$ & $3^{4}=81$\\\hline
une fraction & la touche \fbox{\small$\div$} crée une fraction & $\frac34+\frac25=\frac{23}{20}$\\\hline
reprendre un calcul & flèche \fbox{\small$\uparrow$}, puis \fbox{\small OK} & le calcul revient dans la ligne de saisie\\\hline
effacer l'historique & \fbox{\small shift} puis \fbox{\small$\leftarrow$} (\texttt{clear}) & écran vide\\\hline
une puissance de 10 & touche \fbox{\small$\times10^{x}$} : $3\times10^{2}+1{,}5\times10^{3}$ & $1~800$\\\hline
décomposer un entier & \fbox{\small toolbox}, \texttt{Arithmétique}, \texttt{factor(360)} & $2^{3}\times3^{2}\times5$\\\hline
le nombre $\pi$ & \fbox{\small shift} \fbox{\small$\times10^{x}$} : $2\times\pi\times 3$ & $\approx 18{,}85$\\\hline
\end{tabular}
\end{center}
Un nombre négatif se tape entre parenthèses : $2\times(-3)^{2}$ vaut $18$, mais
$2\times -3^{2}$ est compris comme $2\times(-9)=-18$. {\footnotesize\itshape Les intitulés
peuvent varier selon la version du logiciel.}
\end{proprietebox}

\begin{proprietebox}[Application \texttt{Grapheur}]
\begin{center}
\renewcommand{\arraystretch}{1.3}
\begin{tabular}{|>{\raggedright\arraybackslash}p{4.3cm}|>{\raggedright\arraybackslash}p{12.3cm}|}
\hline
\rowcolor{coulProp!10}\textbf{Pour} & \textbf{On fait}\\\hline
entrer une fonction & Onglet \texttt{Expressions}, \texttt{Ajouter un élément}, puis
  par exemple $f(x)=x^{2}-2$ ($x$ avec la touche \fbox{\small$x,n,t$}) et
  \fbox{\small OK}.\\\hline
afficher un tableau & Onglet \texttt{Tableau}, \texttt{Régler l'intervalle} :
  \texttt{X début}, \texttt{X fin}, \texttt{Pas}.\\\hline
tracer la courbe & Onglet \texttt{Graphique}, menu \texttt{Axes} pour régler $X_{\min}$,
  $X_{\max}$, $Y_{\min}$, $Y_{\max}$.\\\hline
lire une image & Sur la courbe : \fbox{\small OK}, \texttt{Aller à}, puis la valeur de
  $x$.\\\hline
chercher des antécédents & Sur la courbe : \fbox{\small OK}, \texttt{Calculer},
  \texttt{Antécédent}, puis la valeur ; les flèches passent d'un antécédent à
  l'autre.\\\hline
couper deux courbes & Sur l'une des courbes : \fbox{\small OK}, \texttt{Calculer},
  \texttt{Intersection}.\\\hline
\end{tabular}
\end{center}
\end{proprietebox}

\begin{methodebox}[Méthode 13 --- Étudier une fonction à la calculatrice]
On entre la fonction, on règle le tableau (début, fin, pas), puis les axes d'après le
tableau ; on lit enfin images et antécédents.
\etapeexemples
Soit ${f(x)=x^{2}-2}$.
\begin{questions}[label=\alph*)]
  \item Afficher le tableau de valeurs de $f$ de $-4$ à $4$ avec un pas de $1$.
  \item Tracer la courbe de $f$ avec des axes adaptés.
  \item Lire $f(3)$ et $f(1{,}5)$, puis les antécédents de $3$ arrondis à $0{,}01$.
\end{questions}
\begin{correction}
\cas{a)}
\explique{\texttt{Régler l'intervalle} : \texttt{X début} $-4$, \texttt{X fin} $4$,
\texttt{Pas} $1$.}{début, fin et pas}
\begin{center}
\footnotesize
\renewcommand{\arraystretch}{1.2}
\begin{tabular}{|c|*{9}{c|}}
\hline
$x$ & $-4$ & $-3$ & $-2$ & $-1$ & $0$ & $1$ & $2$ & $3$ & $4$\\\hline
$f(x)$ & $14$ & $7$ & $2$ & $-1$ & $-2$ & $-1$ & $2$ & $7$ & $14$\\\hline
\end{tabular}
\end{center}
\cas{b)}
\explique{\texttt{Axes} : $X_{\min}=-4$, $X_{\max}=4$, $Y_{\min}=-3$, $Y_{\max}=15$.}{un
peu au-delà de la plus petite et de la plus grande image}
\cas{c)}
\explique{\texttt{Aller à} : $f(3)=7$ et $f(1{,}5)=0{,}25$.}{lecture d'une image}
\explique{\texttt{Calculer}, \texttt{Antécédent}, valeur $3$ : $x\approx -2{,}24$ et
$x\approx 2{,}24$.}{deux antécédents, donnés arrondis}
\end{correction}

\etape{À vous de jouer}
Soit ${g(x)=(x+1)^{2}}$.
\begin{questions}
  \item Afficher le tableau de valeurs de $g$ de $-3$ à $3$ avec un pas de $0{,}5$, et le
        recopier.
  \item Tracer la courbe de $g$ avec des axes adaptés.
  \item Lire $g(0{,}25)$, puis les antécédents de $3$ arrondis à $0{,}01$.
\end{questions}
\reponse{1\textdegree) $4$ ; $2{,}25$ ; $1$ ; $0{,}25$ ;
$0$ ; $0{,}25$ ; $1$ ; $2{,}25$ ; $4$ ; $6{,}25$ ; $9$ ; $12{,}25$ ; $16$.\;
2\textdegree) par exemple $x$ de $-3$ à $3$, $y$ de $-1$ à $17$.\;
3\textdegree) $g(0{,}25)=1{,}5625$ ;\; antécédents de $3$ : $\approx -2{,}73$ et
$\approx 0{,}73$.}
\end{methodebox}

\afaire{exercices 29 à 32}

\partiecours{Résolution graphique}

\begin{proprietebox}[Propriété 3 --- Résoudre graphiquement]
\begin{itemize}
  \item Les solutions de $f(x)=k$ sont les \textbf{abscisses} des points de $\Cf$ situés
        sur la droite horizontale $y=k$ : ce sont les antécédents de $k$.
  \item Les solutions de $f(x)>k$ sont les abscisses des points de $\Cf$ situés
        \textbf{strictement au-dessus} de cette droite ; pour $f(x)<k$,
        \textbf{strictement au-dessous}.
  \item Pour $f(x)\ge k$ ou $f(x)\le k$, on garde aussi les points sur la droite :
        crochets fermés.
\end{itemize}
Une équation se conclut par des \textbf{accolades}, une inéquation par des
\textbf{intervalles}.
\end{proprietebox}

\begin{methodebox}[Méthode 14 --- Résoudre graphiquement une équation ou une inéquation]
\etapeexemples
Un rectangle a un périmètre de $12$~cm ; un de ses côtés mesure $x$~cm, avec
$0\le x\le 6$. Son aire est ${f(x)=6x-x^{2}}$.
\begin{questions}[label=\alph*)]
  \item Dresser le tableau de valeurs de $f$ de $0$ à $6$ avec un pas de $1$, puis tracer
        la courbe de $f$.
  \item Résoudre graphiquement $f(x)=5$, puis $f(x)=10$.
  \item Résoudre graphiquement $f(x)\ge 5$, puis $f(x)<5$.
\end{questions}
\begin{correction}
\begin{minipage}[t]{0.58\linewidth}\raggedright
\cas{a)}\par\nopagebreak
\begin{center}
\renewcommand{\arraystretch}{1.25}
\begin{tabular}{|c|*{7}{c|}}
\hline
$x$ & $0$ & $1$ & $2$ & $3$ & $4$ & $5$ & $6$\\\hline
$f(x)$ & $0$ & $5$ & $8$ & $9$ & $8$ & $5$ & $0$\\\hline
\end{tabular}
\end{center}
\cas{b)}\par\nopagebreak
La droite $y=5$ coupe $\Cf$ aux points d'abscisses $1$ et $5$ : ${S=\ens{1\,;5}}$.
La droite $y=10$ ne coupe pas $\Cf$ : ${S=\varnothing}$.

\cas{c)}\par\nopagebreak
$\Cf$ est sur la droite $y=5$ ou au-dessus entre $1$ et $5$ : ${S=\intff{1}{5}}$
(en rouge sur l'axe).
$\Cf$ est strictement au-dessous au début et à la fin :
${S=\intfo{0}{1}\cup\intof{5}{6}}$.
\end{minipage}\hfill
\begin{minipage}[t]{0.38\linewidth}
\vspace{0pt}\centering
\begin{tikzpicture}[x=0.42cm,y=0.42cm]
  \repere{-1}{-1}{7}{10}
  \gradx{1,2,3,4,5,6}
  \grady{1,2,3,4,5,6,7,8,9}
  \draw[coulCourbe,line width=1.1pt,domain=0:6,samples=60] plot (\x,{6*(\x)-(\x)*(\x)});
  \foreach \a/\b in {0/0,1/5,2/8,3/9,4/8,5/5,6/0} \path[coulCourbe] (\a,\b) node[croix]{};
  \node[coulCourbe,font=\footnotesize] at (3,9.7) {$\Cf$};
  \draw[black!70,thick] (-1,5)--(7,5);
  \node[black!70,font=\scriptsize,above] at (6.3,5) {$y=5$};
  \draw[black!70,dashed] (1,5)--(1,0) (5,5)--(5,0);
  \draw[red,line width=1.6pt] (1,0)--(5,0);
\end{tikzpicture}
\end{minipage}
\explique{Le rectangle a une aire d'au moins $5$~cm$^{2}$ quand son côté mesure entre
$1$ et $5$~cm.}{on revient au problème}
\end{correction}

\begin{remarquebox}
Une lecture graphique donne en général des valeurs \textbf{approchées} ; ici, le calcul
confirme les valeurs lues : $f(1)=6-1=5$ et $f(5)=30-25=5$. On ne conclut que sur la
partie tracée de la courbe.
\end{remarquebox}

\etape{À vous de jouer}
\begin{minipage}[t]{0.56\linewidth}\raggedright
Ci-contre, la courbe d'une fonction $f$ définie sur $\intff{-4}{4}$.
\begin{questions}
  \item Résoudre graphiquement : \quad $f(x)=1$ \quad $f(x)=2$ \quad $f(x)=-2$ \quad
        $f(x)=3$
  \item Résoudre graphiquement : \quad $f(x)\ge 1$ \quad $f(x)<2$ \quad $f(x)>-2$
\end{questions}
\end{minipage}\hfill
\begin{minipage}[t]{0.40\linewidth}
\vspace{0pt}\centering
\begin{tikzpicture}[x=0.45cm,y=0.45cm]
  \repere{-5}{-3}{5}{3}
  \gradx{-4,-3,-2,-1,1,2,3,4}
  \grady{-2,-1,1,2}
  \draw[coulCourbe,line width=1.1pt] (-4,-1)
    .. controls (-3.667,-0.333) and (-3.333,0.5) .. (-3,1)
    .. controls (-2.667,1.5) and (-2.333,2) .. (-2,2)
    .. controls (-1.667,2) and (-1.333,1.5) .. (-1,1)
    .. controls (-0.667,0.5) and (-0.333,-0.5) .. (0,-1)
    .. controls (0.333,-1.5) and (0.667,-2) .. (1,-2)
    .. controls (1.333,-2) and (1.667,-1.5) .. (2,-1)
    .. controls (2.333,-0.5) and (2.667,0.5) .. (3,1)
    .. controls (3.333,1.5) and (3.667,1.667) .. (4,2);
  \path[coulCourbe] (-4,-1) node[croix]{} (4,2) node[croix]{};
  \node[coulCourbe,font=\footnotesize] at (-3.6,2) {$\Cf$};
\end{tikzpicture}
\end{minipage}
\reponse{1\textdegree) $S=\ens{-3\,;-1\,;3}$ ;\; $S=\ens{-2\,;4}$ ;\; $S=\ens{1}$ ;\;
$S=\varnothing$.\; 2\textdegree) $S=\intff{-3}{-1}\cup\intff{3}{4}$ ;\;
$S=\intfo{-4}{-2}\cup\intoo{-2}{4}$ ;\; $S=\intfo{-4}{1}\cup\intof{1}{4}$.}
\end{methodebox}

\afaire{exercices 33 à 36}

\partiecours{Comparer deux fonctions}

\begin{proprietebox}[Propriété 4 --- Deux courbes]
Soit $f$ et $g$ deux fonctions définies sur un même intervalle.
\begin{itemize}
  \item Les solutions de $f(x)=g(x)$ sont les abscisses des \textbf{points
        d'intersection} de $\Cf$ et $\Cg$.
  \item Les solutions de $f(x)<g(x)$ sont les abscisses des points où $\Cf$ est
        \textbf{strictement au-dessous} de $\Cg$ ; pour $f(x)>g(x)$,
        \textbf{strictement au-dessus}.
\end{itemize}
Pour $f(x)\le g(x)$ ou $f(x)\ge g(x)$, on ajoute les abscisses des points
d'intersection.
\end{proprietebox}

\begin{methodebox}[Méthode 15 --- Comparer graphiquement deux fonctions]
\etapeexemples
Sur $\intff{-1}{4}$, on considère ${f(x)=x^{2}-2x}$ et $g(x)=x$, tracées ci-contre.
Résoudre graphiquement :
\begin{questions}[label=\alph*)]
  \item $f(x)=g(x)$
  \item $f(x)<g(x)$
  \item $f(x)\ge g(x)$
\end{questions}
\begin{correction}
\begin{minipage}[t]{0.58\linewidth}\raggedright
\cas{a)}\par\nopagebreak
Les courbes se coupent aux points d'abscisses $0$ et $3$ : ${S=\ens{0\,;3}}$.

\cas{b)}\par\nopagebreak
$\Cf$ est strictement au-dessous de $\Cg$ entre $0$ et $3$, bornes exclues :
${S=\intoo{0}{3}}$ (en rouge sur l'axe).

\cas{c)}\par\nopagebreak
$\Cf$ est au-dessus de $\Cg$, ou la coupe, sur le reste de l'intervalle :
${S=\intff{-1}{0}\cup\intff{3}{4}}$.

\smallskip
\emph{Vérification : $f(0)=0$ et $g(0)=0$ ; $f(3)=9-6=3$ et $g(3)=3$.}
\end{minipage}\hfill
\begin{minipage}[t]{0.38\linewidth}
\vspace{0pt}\centering
\begin{tikzpicture}[x=0.45cm,y=0.45cm]
  \repere{-2}{-2}{5}{9}
  \gradx{-1,1,2,3,4}
  \grady{-1,1,2,3,4,5,6,7,8}
  \draw[coulCourbe,line width=1.1pt,domain=-1:4,samples=60] plot (\x,{(\x)*(\x)-2*(\x)});
  \draw[coulCourbeC,line width=1.1pt] (-1,-1)--(4,4);
  \path (0,0) node[croix]{} (3,3) node[croix]{};
  \draw[red,line width=1.6pt] (0,0)--(3,0);
  \draw[dashed] (3,3)--(3,0);
  \node[coulCourbe,font=\footnotesize,left] at (3.9,8) {$\Cf$};
  \node[coulCourbeC,font=\footnotesize,right] at (4,4) {$\Cg$};
\end{tikzpicture}
\end{minipage}
\end{correction}

\etape{À vous de jouer}
\begin{minipage}[t]{0.56\linewidth}\raggedright
Sur $\intff{-3}{2}$, on considère ${f(x)=4-x^{2}}$ et $g(x)=x+2$, tracées ci-contre.
Résoudre graphiquement :
\begin{questions}
  \item $f(x)=g(x)$
  \item $f(x)\le g(x)$
  \item $f(x)>g(x)$
\end{questions}
\reponse{1\textdegree) $S=\ens{-2\,;1}$ ;\; 2\textdegree) $S=\intff{-3}{-2}\cup\intff{1}{2}$ ;\;
3\textdegree) $S=\intoo{-2}{1}$.}
\end{minipage}\hfill
\begin{minipage}[t]{0.40\linewidth}
\vspace{0pt}\centering
\begin{tikzpicture}[x=0.45cm,y=0.45cm]
  \repere{-4}{-6}{3}{5}
  \gradx{-3,-2,-1,1,2}
  \grady{-5,-4,-3,-2,-1,1,2,3,4}
  \draw[coulCourbe,line width=1.1pt,domain=-3:2,samples=60] plot (\x,{4-(\x)*(\x)});
  \draw[coulCourbeC,line width=1.1pt] (-3,-1)--(2,4);
  \node[coulCourbe,font=\footnotesize,right] at (-3,-5) {$\Cf$};
  \node[coulCourbeC,font=\footnotesize,left] at (2,4.4) {$\Cg$};
\end{tikzpicture}
\end{minipage}
\end{methodebox}

\afaire{exercices 37 à 39}

\end{document}
